\documentclass[11pt]{article}

\usepackage[margin=1in]{geometry}
\usepackage[colorlinks=true, linkcolor=blue, urlcolor=blue]{hyperref}
\usepackage{enumitem}
\usepackage{parskip}

\title{CalifAI: How Everything Works\\
\large A Chrome Extension for Capturing Calendar Events with AI}
\author{Adrian Hito}
\date{June 2026}

\begin{document}

\maketitle
\tableofcontents
\newpage

% ===================================================================
\section{What the Extension Does}

CalifAI is a Chrome extension (Manifest V3) that turns anything visible on a
webpage into Google Calendar events in three steps:

\begin{enumerate}
  \item \textbf{Capture}: the user screenshots either the full visible tab or a
        dragged selection of the page.
  \item \textbf{Extract}: the screenshot is sent to a vision-capable AI model
        (Google Gemini 2.5 Flash by default, OpenAI GPT-4o-mini as an
        alternative) which returns structured JSON describing every calendar
        event it can find: title, start/end times, location, description,
        recurrence rules, and all-day status.
  \item \textbf{Import}: the user reviews and edits the extracted events, then
        the extension creates them in Google Calendar via the Calendar REST
        API using OAuth.
\end{enumerate}

% ===================================================================
\section{Tech Stack and Project Layout}

\begin{itemize}
  \item \textbf{WXT} -- the extension build framework (wraps Vite). Configured
        in \texttt{wxt.config.ts}, which also defines the manifest
        (permissions, OAuth client ID, scopes). Output goes to
        \texttt{.output/chrome-mv3/}.
  \item \textbf{React + TypeScript} -- the popup and options page UIs.
  \item \textbf{Zustand} -- popup state management
        (\texttt{hooks/useAppState.ts}), persisted to
        \texttt{chrome.storage.session} on every mutation.
  \item \textbf{Zod} -- runtime validation of AI responses
        (\texttt{lib/ai/schema.ts}).
\end{itemize}

The extension has three execution contexts:

\begin{description}
  \item[Popup] (\texttt{entrypoints/popup/}) -- the React app the user sees.
        It dies whenever it loses focus, which drives much of the
        architecture.
  \item[Background service worker] (\texttt{entrypoints/background/}) -- does
        everything privileged: screen capture, AI calls, OAuth, Calendar API.
        MV3 service workers are killed aggressively by Chrome when idle.
  \item[Content script] (\texttt{entrypoints/content.ts}) -- injected into the
        page to draw the drag-to-select overlay and report the selected
        rectangle back.
\end{description}

These three communicate via \texttt{chrome.runtime.sendMessage} with a typed
protocol (\texttt{lib/messaging/types.ts}). The popup sends messages such as
\texttt{CAPTURE\_AND\_EXTRACT}, \texttt{CREATE\_EVENTS},
\texttt{AUTHORIZE\_GOOGLE}; the service worker routes them in
\texttt{lib/messaging/handlers.ts}.

% ===================================================================
\section{Popup State Machine}

The popup is a view router (\texttt{entrypoints/popup/App.tsx}) switching on a
single \texttt{view} field:
\texttt{setup} $\rightarrow$ \texttt{home} $\rightarrow$ \texttt{loading}
$\rightarrow$ \texttt{event-selection} / \texttt{review} $\rightarrow$
\texttt{edit} $\rightarrow$ \texttt{importing} $\rightarrow$ \texttt{success},
with \texttt{error} reachable from anywhere.

Key state fields (\texttt{types/app-state.ts}): \texttt{events} (the working
list), \texttt{currentEvent}, \texttt{selectedEventIndices} (multi-select),
\texttt{viewHistory} (a stack enabling \texttt{goBack()}), \texttt{error}
(interrupting), and \texttt{notice} (a non-interrupting banner, e.g.\ ``Added
3 events'').

\textbf{Persistence and hydration.} Every state change is saved to
\texttt{chrome.storage.session} under the key \texttt{appState}. When the
popup opens, \texttt{hydrate()} restores it. Crucially, \texttt{App.tsx}
renders \emph{nothing} until hydration completes: if it rendered the default
state first, a phantom \texttt{HomeView} would mount, start its own
capture-result handler, and race the real view's handler (this once caused a
bug where the home handler \emph{replaced} the user's event list while the
real handler then saw every new event as a duplicate).

\textbf{Navigation rules.} \texttt{setView} pushes the previous view onto
\texttt{viewHistory} unless the previous view was \texttt{loading} or
\texttt{edit} (transient views never become ``back'' targets). The Cancel
button on the loading view therefore calls \texttt{goBack()} -- returning to
home for initial captures, or to review/event-selection for ``Capture More''
-- after clearing the pending-capture flags.

% ===================================================================
\section{The Capture Pipeline}

\subsection{Full-screen capture}

Simple, synchronous round trip: popup sends \texttt{CAPTURE\_AND\_EXTRACT}
with \texttt{useSelection: false}; the worker calls
\texttt{chrome.tabs.captureVisibleTab} (JPEG, quality 85), runs extraction,
and returns the events in the message response. The popup stays open the
whole time showing the loading view.

\subsection{Area selection capture (the interesting one)}

Chrome has \textbf{no API to capture a sub-region} of a page --
\texttt{captureVisibleTab} always returns the full viewport. Additionally,
the popup \textbf{closes itself} so the user can interact with the page. The
flow:

\begin{enumerate}
  \item Popup clears stale capture state, records
        \texttt{lastCaptureMode = 'area'} in session storage, fires the
        message, and calls \texttt{window.close()}.
  \item The worker pings the content script (injecting it via
        \texttt{chrome.scripting} if absent) and sends
        \texttt{START\_SELECTION}.
  \item The content script draws a full-screen overlay with a crosshair
        cursor and a selection box (the dark surround is a huge
        \texttt{box-shadow}). On mouse-up it returns the rectangle in CSS
        pixels plus \texttt{window.devicePixelRatio}. Selections under
        $10\times10$ px are rejected as accidental clicks (originally
        $50\times50$, which wrongly rejected single event rows).
  \item The worker captures the full viewport, then \emph{immediately} crops
        to the selection in an \texttt{OffscreenCanvas}
        (\texttt{lib/utils/image-crop.ts}). CSS coordinates are multiplied by
        \texttt{devicePixelRatio}, which in Chrome already incorporates page
        zoom, so the mapping to bitmap pixels is exact. The crop and the
        downscale to the provider's target width (768px for Gemini, 1280px
        for OpenAI) happen in \emph{one} canvas pass with a single JPEG
        encode (quality 0.7). The full-frame image exists only transiently in
        worker memory -- it is never stored or transmitted.
  \item The worker writes \texttt{processingCapture: true} and
        \texttt{processingStartedAt} (a timestamp) to session storage,
        reopens the popup via \texttt{chrome.action.openPopup()} (3 retries;
        falls back to a desktop notification), and continues extraction in
        the background.
  \item When done, it writes \texttt{captureComplete: true} plus either
        \texttt{captureResult} or \texttt{captureError}.
\end{enumerate}

\subsection{How the reopened popup picks up the result}

The popup and the worker share a small contract in
\texttt{lib/storage/capture-session.ts}:

\begin{itemize}
  \item \texttt{waitForCaptureCompletion()} resolves the moment the result
        lands, by listening to \texttt{chrome.storage.session.onChanged}
        (instant) with a 2-second fallback poll (in case an event is missed)
        and a 2-minute timeout. It resolves with one of three statuses:
        \texttt{complete}, \texttt{timeout}, or \texttt{cancelled} (the
        processing flag disappeared without a result, i.e.\ the user pressed
        Cancel).
  \item \textbf{Staleness protection}: a \texttt{processingCapture} flag older
        than 2 minutes (per \texttt{processingStartedAt}) is treated as
        garbage left by a killed service worker. Without this, one crashed
        extraction permanently trapped the popup in an infinite loading loop
        (Cancel went home, home saw the flag, bounced back to loading).
  \item Cancel on the loading view calls \texttt{clearCaptureSession()}
        \emph{before} navigating, so no view re-detects the flag.
\end{itemize}

Who waits depends on where the user was: \texttt{HomeView} handles initial
captures (it \emph{replaces} the event list), while the shared
\texttt{useCaptureMore} hook (mounted by both ReviewView and
EventSelectionView) handles ``Capture More'' (it \emph{merges}).

\subsection{Capture More and merging}

``Capture More'' reuses whatever mode the user originally chose
(\texttt{lastCaptureMode}). Results are merged into the existing list with
duplicate detection keyed on \texttt{title + startDate}. Outcomes
\emph{never} route to the error view (that would strand the user away from
their events); instead the user always lands on the event list with a
\texttt{notice} banner: ``Added $X$ events ($Y$ duplicates skipped)'',
``Everything you captured is already on the list'', or ``No new events found
in that capture''. The banner auto-dismisses after 5 seconds.

% ===================================================================
\section{AI Extraction}

\subsection{Provider abstraction}

\texttt{lib/ai/provider-registry.ts} maps a name to an \texttt{AIProvider}
implementation with one method: \texttt{extractEvents(imageBase64, apiKey)}.
Two providers exist: Gemini (default, free tier) and OpenAI.
\texttt{lib/ai/extraction.ts} orchestrates: throttle (min 1s between
requests), read settings, resize the image if it is still wider than the
provider target (skipped entirely when the crop already downscaled it), and
dispatch.

\subsection{The prompt}

\texttt{lib/ai/extraction-prompt.ts} exports
\texttt{getExtractionPrompt()}. Design points:

\begin{itemize}
  \item Detailed rules: extract \emph{every} event; never hallucinate; if an
        event has a start time but no end time, add exactly one hour; date
        only $\Rightarrow$ all-day event; detect recurrence keywords
        (``every Monday'', ``bi-weekly'', ``8 sessions'') and map them to
        frequency/interval/count/until/byDay.
  \item The long instruction body is \textbf{byte-identical across requests},
        and the current date (day granularity plus weekday, e.g.\ ``Current
        date: Wednesday, 2026-06-10'') is appended at the \emph{end}. A
        stable prefix lets Gemini's implicit prompt caching discount and
        speed up repeat requests; the original code embedded a
        millisecond-precision timestamp at the \emph{start}, which defeated
        caching entirely.
\end{itemize}

\subsection{Gemini specifics (\texttt{lib/ai/gemini-provider.ts})}

\begin{itemize}
  \item Model \texttt{gemini-2.5-flash} via the
        \texttt{generativelanguage.googleapis.com} REST endpoint.
  \item The API key is sent in the \texttt{x-goog-api-key} \textbf{header},
        not the legacy \texttt{?key=} query parameter. This matters because
        Google's newer \texttt{AQ.}-prefix ``auth keys'' are misparsed as
        OAuth tokens when sent as a query parameter, producing the confusing
        401 ``Request had invalid authentication credentials. Expected OAuth
        2 access token\ldots'' (\texttt{ACCESS\_TOKEN\_TYPE\_UNSUPPORTED}).
        That error looks like a Google-account problem but is an API-key
        transport problem.
  \item \texttt{thinkingConfig: \{ thinkingBudget: 0 \}} -- Gemini 2.5 Flash
        runs hidden ``dynamic thinking'' by default, which adds large latency
        and burns output tokens. Structured extraction with a strict schema
        doesn't need it; disabling it is the single biggest speed win and
        also reduces token consumption.
  \item \texttt{responseMimeType: application/json} plus a strict
        \texttt{responseSchema} forces valid JSON matching the event shape.
        The schema deliberately has \emph{no} \texttt{reasoning} field --
        generating an explanation paragraph nobody reads wasted output tokens
        (the slowest part of any LLM response) on every capture.
  \item Retries: up to 3 attempts with exponential backoff (2s/4s/8s), but
        only for transient errors (503 ``high demand'', quota). Error
        mapping translates Google's raw messages into actionable ones:
        quota $\rightarrow$ ``check your limits''; 400/401/403 key problems
        $\rightarrow$ ``Gemini rejected your API key'' (deliberately
        \emph{not} echoing the OAuth-sounding raw text); truncated JSON
        $\rightarrow$ ``capture fewer events at once''.
\end{itemize}

\subsection{Validation}

Responses are parsed with Zod (\texttt{AIExtractionResponseSchema}). A
transform scrubs junk the model sometimes emits in text fields (e.g.\
\texttt{location: "(no location given)"} becomes \texttt{undefined}).

% ===================================================================
\section{Google OAuth}

\subsection{Why \texttt{launchWebAuthFlow} and not \texttt{getAuthToken}}

\texttt{chrome.identity.getAuthToken} is pinned to the Chrome profile's
signed-in account: it can \emph{never} show Google's account chooser, so
``switch account'' was impossible, and Chrome caches tokens that Google has
already revoked (causing mysterious 401s that a naive cache check cannot
detect). The extension therefore uses
\texttt{chrome.identity.launchWebAuthFlow}
(\texttt{entrypoints/background/google-auth.ts}):

\begin{itemize}
  \item \textbf{Implicit grant}: the auth URL requests
        \texttt{response\_type=token}; the access token comes back in the
        redirect URL fragment. No client secret is involved (and none should
        ever ship in an extension).
  \item \textbf{Redirect URI}: \texttt{https://<extension-id>.chromiumapp.org/}
        -- a virtual address Chrome intercepts. It must be registered as an
        authorized redirect URI on a \emph{Web application} OAuth client in
        Google Cloud Console (extension-type clients reject it).
  \item \textbf{Connect / Switch account}: always an interactive flow with
        \texttt{prompt=select\_account}, which lands directly on Google's
        ``Choose an account'' page. The existing grant is deliberately
        \emph{not} revoked first -- revoking forces full re-consent, which
        shows the scary ``Google hasn't verified this app'' page (the consent
        screen is in testing mode, so only listed test users can authorize,
        and first-time consent always shows the unverified warning).
  \item \textbf{Token storage}: access token and expiry live in
        \texttt{chrome.storage.local} (\texttt{googleAccessToken},
        \texttt{googleTokenExpiresAt}). Tokens are refreshed silently
        (\texttt{interactive: false}, \texttt{prompt=none}) using a stored
        \texttt{login\_hint} (the connected account's email) so Google knows
        which session to mint for without UI.
  \item \textbf{Cancel safety}: closing the chooser rejects the flow; the
        stored token is only replaced after \emph{successful} auth, so the
        current account stays connected and displayed.
  \item \texttt{AuthRequiredError} (code \texttt{AUTH\_REQUIRED}) is thrown
        when reconnection is genuinely required; the error view recognizes it
        (and auth-flavored message text) and shows a \textbf{Connect Google
        Account} button instead of a useless ``Try Again''.
\end{itemize}

\subsection{Calendar API (\texttt{google-calendar.ts})}

\texttt{makeAuthenticatedRequest} attaches \texttt{Authorization: Bearer
<token>}. On a 401 it discards the stored token and retries once with a
freshly minted one; if that also 401s, it throws
\texttt{AuthRequiredError} with a human message rather than leaking Google's
raw error. Event conversion maps the internal event shape to the Calendar v3
body: all-day events use \texttt{date} (not \texttt{dateTime}); timed events
attach the local IANA timezone; recurrence becomes an RFC~5545
\texttt{RRULE} string (\texttt{FREQ}, \texttt{INTERVAL}, \texttt{COUNT},
\texttt{UNTIL}, \texttt{BYDAY}); reminders become
\texttt{reminders.overrides}; an optional \texttt{colorId} (1--11) colors the
event.

% ===================================================================
\section{Error-Handling Philosophy}

\begin{itemize}
  \item \textbf{Classify, then act.} The error view inspects the message/code
        and offers the action that can actually fix it: auth errors
        $\rightarrow$ Connect Google Account; API-key and rate-limit errors
        $\rightarrow$ Change API Keys (opens settings); selection errors
        (too small / cancelled) $\rightarrow$ Capture Again \emph{relaunches
        the selection overlay directly} (it resets the persisted view to
        home first so the reopened popup picks up the pending result);
        everything else $\rightarrow$ Try Again / Capture Again.
  \item \textbf{Never strand the user.} Merge outcomes use the
        \texttt{notice} banner instead of the error view; Cancel goes
        \emph{back}, not home; cancelling auth keeps the current account.
  \item \textbf{Never trust a flag without a timestamp.} Anything a killable
        service worker writes (\texttt{processingCapture}) carries
        \texttt{processingStartedAt} and is ignored when stale.
\end{itemize}

% ===================================================================
\section{Performance Decisions (and Why)}

\begin{enumerate}
  \item \textbf{Thinking disabled} (\texttt{thinkingBudget: 0}): the largest
        latency reduction, also cheaper.
  \item \textbf{Fewer output tokens}: no \texttt{reasoning} field; output
        tokens dominate response time.
  \item \textbf{Cache-friendly prompt}: stable prefix, day-granularity date
        at the end.
  \item \textbf{Single-pass image pipeline}: crop + downscale + encode once
        (was: three JPEG encode/decode generations, which also blurred the
        text the model must read).
  \item \textbf{No dead storage writes}: the full captured image was once
        written to session storage and never read -- a megabyte-class
        serialization on the critical path, removed.
  \item \textbf{Event-driven results}: \texttt{storage.onChanged} instead of
        500\,ms polling (average 250\,ms saved after extraction finishes).
  \item \textbf{Image sizing}: Gemini tiles images at $768\times768$; sending
        768px-wide images minimizes vision tokens.
\end{enumerate}

% ===================================================================
\section{Likely Quiz Questions}

\begin{description}
  \item[Why does the popup close during area selection?] Chrome popups cannot
        coexist with page interaction; the overlay needs mouse events on the
        page. State is handed off through \texttt{chrome.storage.session}
        because the popup's JS context is destroyed.
  \item[Why capture the whole screen and crop?] There is no region-capture
        API in Chrome extensions. The full frame lives only in worker memory
        for milliseconds; only the crop is stored/sent.
  \item[Why did ``invalid authentication credentials'' appear on capture?]
        The Gemini API key (new \texttt{AQ.} format) was sent as
        \texttt{?key=} and misparsed as an OAuth token. Fix: send keys via
        \texttt{x-goog-api-key} header.
  \item[Why \texttt{launchWebAuthFlow}?] \texttt{getAuthToken} cannot show an
        account chooser and caches revoked tokens. The web flow gives full
        control (\texttt{prompt=select\_account}) at the cost of managing
        token storage/refresh manually.
  \item[How are duplicates detected when capturing more?] Set membership on
        the composite key \texttt{title + startDate}.
  \item[What stops an infinite loading loop if the worker dies?] The
        timestamped processing flag: readers treat flags older than 2
        minutes as failed, and Cancel clears the flags before navigating.
  \item[Why render nothing before hydration?] To prevent a phantom HomeView
        from racing the real view's capture handler and replacing the event
        list.
  \item[Where do API keys and tokens live?] \texttt{chrome.storage.local}
        (persistent): AI keys, Google token, settings.
        \texttt{chrome.storage.session} (cleared with the browser): app
        state, capture handoff flags. Nothing sensitive is hardcoded.
\end{description}

\end{document}
